%%%% ADT surface figure %%%%%
\begin{comment}
\begin{table}[ht]
\centering
\begin{tabular}{ccc} % Adjust for three columns
\small
GT Mesh & GT Depth Fusion & ZoeDepth \\ % Headers for the first row of images
\includegraphics[width=0.3\textwidth]{figures/adt_mesh_viz/gt_mesh.jpg} & 
\includegraphics[width=0.3\textwidth]{figures/adt_mesh_viz/gt_fusion.jpg} & 
\includegraphics[width=0.3\textwidth]{figures/adt_mesh_viz/zoe.jpg} \\ 
SimpleRecon & NeuralRecon & \OCCmodel{} \\ % Headers for the second row of images
\includegraphics[width=0.3\textwidth]{figures/adt_mesh_viz/simplerecon.jpg} & 
\includegraphics[width=0.3\textwidth]{figures/adt_mesh_viz/neural_recon.jpg} & 
\includegraphics[width=0.3\textwidth]{figures/adt_mesh_viz/efm.jpg} \\
\end{tabular}
\caption{Surface reconstruction on the ADT work 106 sequence. For ZoeDepth and SimpleRecon we use the pretrained models while for \OCCmodel{} and NeuralRecon we re-trained on the ASE 1k training set. \TODO{Add NeuralRecon; Check Zoe isse; Add zoom-in view; GT Mesh and GT Depth shading.} }
\label{fig:mesh_adt}
\end{table}
\end{comment}


%%%% ASE surface figure %%%%%
\begin{figure}[htbp]
    \centering
    \begin{subfigure}[b]{\subfigwidth\textwidth}
        \centering
        \caption{\tiny \textbf{GT Mesh}}
        \includegraphics[width=\textwidth]{figures/adt_mesh_viz/gt_mesh_comb.png}
    \end{subfigure}
    \hfill
    \begin{subfigure}[b]{\subfigwidth \textwidth}
        \centering
        \caption{\tiny \textbf{GT Depth Fusion}}
        \includegraphics[width=\textwidth]{figures/adt_mesh_viz/gt_depth_comb.png}
    \end{subfigure}
    \hfill
    \begin{subfigure}[b]{\subfigwidth\textwidth}
        \centering
        \caption{\tiny \textbf{\EFMmodel{}}}
        \includegraphics[width=\textwidth]{figures/adt_mesh_viz/openocc_comb.png}
    \end{subfigure}
    \hfill
    \begin{subfigure}[b]{\subfigwidth\textwidth}
        \centering
        \caption{\tiny \textbf{NeuralRecon}}
        \includegraphics[width=\textwidth]{figures/adt_mesh_viz/neuralrecon_trained_comb.png}
    \end{subfigure}
    \hfill
    \begin{subfigure}[b]{\subfigwidth\textwidth}
        \centering
        \caption{\tiny \textbf{NeuralRecon OTS}}
        \includegraphics[width=\textwidth]{figures/adt_mesh_viz/neuralrecon_ots_comb.png}
    \end{subfigure}
    \hfill
    \begin{subfigure}[b]{\subfigwidth \textwidth}
        \centering
        \caption{\tiny \textbf{ZoeDepth}}
        \includegraphics[width=\textwidth]{figures/adt_mesh_viz/zoe_comb.png}
    \end{subfigure}
    \hfill
    \begin{subfigure}[b]{\subfigwidth\textwidth}
        \centering
        \caption{\tiny \textbf{SimpleRecon}}
        \includegraphics[width=\textwidth]{figures/adt_mesh_viz/simplerecon_scaling_comb.png}
    \end{subfigure}
    \hfill
    \begin{subfigure}[b]{\subfigwidth\textwidth}
        \centering
        \caption{\tiny \textbf{ConsistentDepth}}
        \includegraphics[width=\textwidth]{figures/adt_mesh_viz/consistent_depth_comb.png}
    \end{subfigure}
    
    
    \caption{Surface reconstruction on ADT-106. It shows how models trained on simulation generalize to real data. \EFMmodel{} has the best generalization ability among all the methods.}
    \label{fig:mesh_adt}
    \vspace{-5mm}
\end{figure}